% \iffalse meta-comment
%
% hit-title.dtx 是各个类共用的论文标题解析
%
% \fi
%
% \section{共享标题解析}
%
% 标题分别保存两段内容及独立的副题目，供固定题目栏、自由排版和书签取用。
% 新的公开输入由 \file{hit-bilingual.dtx} 解析：title 的尾随方括号是副语言，
% subtitle 单独保存副题目。下方解析函数用于内部默认值和历史字段，不能作为新接口直接调用。
%
%    \begin{macrocode}
%<*shared-title>
\tl_new:N \l__hit_title_language_tl
\tl_new:N \l__hit_title_source_tl
%    \end{macrocode}
%
% \begin{macro}{\hit@title@set@first,\hit@title@set@second,\hit@title@set@subtitle}
% 三格各自的写入口。旧键 \option{ctitleone}、\option{ctitletwo}、\option{csubtitle}
% 在兼容期里各写一格，互不覆盖。
%    \begin{macrocode}
\cs_new_protected:Npn \hit@title@set@first #1#2
  {\__hit_info_set:nn {title@first@#1} {#2}}
\cs_new_protected:Npn \hit@title@set@second #1#2
  {\__hit_info_set:nn {title@second@#1} {#2}}
\cs_new_protected:Npn \hit@title@set@subtitle #1#2
  {\__hit_info_set:nn {title@subtitle@#1} {#2}}
%    \end{macrocode}
% \end{macro}
%
% \begin{macro}{\hit@title@parse}
% 解析内部旧式分段题目；匹配后仍有剩余记号时保留整段内容，避免误截以分组开头的普通标题。
%    \begin{macrocode}
\cs_new_protected:Npn \hit@title@parse #1#2
  {
    \tl_set:Nn \l__hit_title_language_tl {#1}
    \tl_set:Nn \l__hit_title_source_tl {#2}
    \hit@title@set@first    {#1} { }
    \hit@title@set@second   {#1} { }
    \hit@title@set@subtitle {#1} { }
    \tl_if_head_is_group:nTF {#2}
      { \hit@title@grab #2 \q_nil \q_stop }
      { \hit@title@set@first {#1} {#2} }
  }
\cs_new_protected:Npn \hit@title@grab #1#2 \q_stop
  {
    \exp_args:NV \hit@title@set@first \l__hit_title_language_tl {#1}
    \__hit_title_after_first:n {#2}
  }
\cs_new_protected:Npn \__hit_title_after_first:n #1
  {
    \tl_if_eq:nnF {#1} { \q_nil }
      {
        \tl_if_head_is_group:nTF {#1}
          { \hit@title@grab@second #1 \q_stop }
          { \__hit_title_after_second:n {#1} }
      }
  }
\cs_new_protected:Npn \hit@title@grab@second #1#2 \q_stop
  {
    \exp_args:NV \hit@title@set@second \l__hit_title_language_tl {#1}
    \__hit_title_after_second:n {#2}
  }
\cs_new_protected:Npn \__hit_title_after_second:n #1
  {
    \tl_if_eq:nnF {#1} { \q_nil }
      {
        \tl_if_head_eq_charcode:nNTF {#1} [
          { \hit@title@grab@subtitle #1 \q_stop }
          { \__hit_title_fallback: }
      }
  }
\cs_new_protected:Npn \hit@title@grab@subtitle [#1] #2 \q_stop
  {
    \exp_args:NV \hit@title@set@subtitle \l__hit_title_language_tl {#1}
    \tl_if_eq:nnF {#2} { \q_nil } { \__hit_title_fallback: }
  }
\cs_new_protected:Npn \__hit_title_fallback:
  {
    \msg_warning:nnV { hithesis } { title-unparsed } \l__hit_title_language_tl
    \exp_args:NVV \hit@title@set@first
      \l__hit_title_language_tl \l__hit_title_source_tl
    \exp_args:NV \hit@title@set@second   \l__hit_title_language_tl { }
    \exp_args:NV \hit@title@set@subtitle \l__hit_title_language_tl { }
  }
%    \end{macrocode}
% \end{macro}
%
% \begin{macro}{\hit@title@cover,\hit@title@flat,\hit@title@line}
% 三种取用。自由排的封面把两段直接接起来，\cs{\textbackslash\textbackslash} 照常换行；
% 声明书与 PDF 书签要一整行，把 \cs{\textbackslash\textbackslash} 屏蔽掉；两条线的封面
% 一次取一段，那种版式里每段都在一个 \cs{makebox} 里，断不了行，同样屏蔽。
% 屏蔽放在分组里，出了这一处不影响别人。
%    \begin{macrocode}
\cs_new:Npn \hit@title@cover #1
  {\tl_use:c {g__hit_info_title@first@#1_tl} \tl_use:c {g__hit_info_title@second@#1_tl}}
\cs_new:Npn \hit@title@flat #1
  {
    \begingroup
      \cs_set:Npn \\ { }
      \tl_use:c {g__hit_info_title@first@#1_tl} \tl_use:c {g__hit_info_title@second@#1_tl}
    \endgroup
  }
\cs_new:Npn \hit@title@line #1#2
  {
    \begingroup
      \cs_set:Npn \\ { }
      \tl_use:c {g__hit_info_title@#2@#1_tl}
    \endgroup
  }
%    \end{macrocode}
%
% 为中英文题目分别生成无参数的读取命令。
%    \begin{macrocode}
\clist_map_inline:Nn \c__hit_lang_clist
  {
    \cs_new:cpn { hit@info@title@cover@#1 }    { \hit@title@cover {#1} }
    \cs_new:cpn { hit@info@title@flat@#1 }     { \hit@title@flat {#1} }
    \cs_new:cpn { hit@info@title@line@one@#1 } { \hit@title@line {#1} { first } }
    \cs_new:cpn { hit@info@title@line@two@#1 } { \hit@title@line {#1} { second } }
  }
%    \end{macrocode}
% \end{macro}
%
% \begin{macro}{\hit@title@warn@wide}
% \changes{v3.2a}{2026/08/07}{题目排不进题目栏时报警告。原先是静默溢出：那一栏是
%   \cs{makebox} 加 \cs{hss}，撑不下就顶出去压着下划线，连 Overfull 都不报}
% 固定题目栏中的 hss 会抑制 Overfull 警告，因此需要单独检测内容宽度。
%    \begin{macrocode}
\box_new:N \l__hit_title_box
\cs_new_protected:Npn \hit@title@warn@wide #1#2#3
  {
    \hbox_set:Nn \l__hit_title_box {#3}
    \dim_compare:nNnT { \box_wd:N \l__hit_title_box } > {#2}
      {
        \msg_warning:nnnxx { hithesis } { title-too-wide } {#1}
          { \dim_to_decimal_in_unit:nn { \box_wd:N \l__hit_title_box } { 1cm } cm }
          { \dim_to_decimal_in_unit:nn {#2} { 1cm } cm }
      }
  }
%    \end{macrocode}
% \end{macro}
%
% \begin{macro}{\hit@define@legacy@title}
% 合并前的四个旧键各写一格：\option{ctitlecover} 与 \option{ctitleone} 写第一段，
% \option{ctitletwo} 写第二段，\option{csubtitle} 与 \option{esubtitle} 写副标题。
% 走通用的 \cs{hit@define@legacy@terms} 不行，那个是整键别名，会互相覆盖。
%    \begin{macrocode}
\cs_new_protected:Npn \hit@define@legacy@title #1#2#3
  {
    \cs_gset_protected_nopar:cpn {#1} ##1
      { \use:c { hit@title@set@#3 } {#2} {##1} }
    \hit@declare@class@key {#1}
    \prop_gput:Nnn \g__hit_renamed_option_prop {#1} { info={title-#2=...} }
  }
%    \end{macrocode}
% \end{macro}
%
% \begin{macro}{\hit@if@subtitle@T}
% show-subtitle 为 auto 时检查副题目是否为空；带语言参数的版本检查指定语种。
%    \begin{macrocode}
\cs_new_protected:Npn \hit@if@subtitle@T #1
  {
    \__hit_tristate_if:NnT \g__hit_subtitle_tl
      { ! \tl_if_empty_p:c {g__hit_info_title@subtitle@zh_tl} } {#1}
  }
\cs_new_protected:Npn \hit@if@subtitle@lang:nT #1#2
  { \__hit_tristate_if:NnT \g__hit_subtitle_tl
      { ! \tl_if_empty_p:c {g__hit_info_title@subtitle@#1_tl} } {#2} }
%    \end{macrocode}
% \end{macro}
%
% \begin{macro}{\hit@define@title}
% 把 \option{title-}\meta{语言} 注册成键，处理器走解析。与 \cs{hit@define@term} 并列，
% 区别只在存值方式：那个原样存一格，这个拆成三格。
%    \begin{macrocode}
\cs_new_protected:Npn \hit@define@title #1
  {
    \hit@declare@class@key { title-#1 }
    \cs_gset_protected_nopar:cpn { title-#1 } ##1 { \hit@title@parse {#1} {##1} }
    \use:c { title-#1 } { }
  }
%    \end{macrocode}
% \end{macro}
%
%    \begin{macrocode}
%</shared-title>
%    \end{macrocode}
%
% \Finale
